\scanpage{161}
\mspage{489}
% position: top continuation
son bon vouloir - par où en somme il nous ``tient'', et (mine de rien) nous manoeu-\linebreak
vre à sa guise...

% position: main paragraph
En bonne logique, l'acceptation d'autrui devrait bien impliquer aussi
l'acceptation de sa façon de voir les choses, qu'elles nous paraissent erronées
ou non, et même lorsqu'il s'agit de sa façon de voir notre propre et précieuse
personne (y inclus nos propres façons de voir...). C'est surtout là pourtant que
le bat blesse - c'est \mbox{l à} le point névralgique dans l'acceptation d'autrui,
et nullement dans l'acceptation de ``défauts'' courants plus ou moins gênants qui
n'impliquent pas directement notre personne. Bien souvent d'ailleurs, si nous \unsure{re|reje-}\linebreak
tons tels ``défauts'' en autrui, c'est surtout parce que par eux nous nous sentons
mis en cause directement, du seul fait d'être confronté à des façons d'être
qui nous paraissent (à tort ou à raison encore) à l'opposé des nôtres. En d'au-\linebreak
tres mots, c'est une \mbox{i n s é c u r i t é} en nous, se manifestant par les
réactions (plus ou moins apparentes ou cachées) de la vanité, qui est le grand
obstacle, s'opposant à notre acceptation d'autrui. Mais cette insécurité pro-\linebreak
fondément enracinée, compensée par les mouvements de la vanité, m'apparaît com-\linebreak
me \unsure{insissolublement|indissolublement} liée à la non-acceptation de nous-même, elle en est comme l'\unsure{o|om-}\linebreak
bre inséparable.

% position: concluding upper paragraph
Ainsi, c'est la pleine acceptation de soi qui apparaît ici comme la clef
qui nous ouvre à l'acceptation d'autrui. Et ce lien qui vient de m'apparaître
ici, rejoint un autre lien profond, que je connais depuis longtemps, depuis tou-\linebreak
jours peut-être: que l'amour de soi est le coeur, paisible et fort, de l'amour
de l'autre.

% position: lower note 111
\(({}^{111})\) (13 octobre) Hier je n'ai pas continué à écrire les notes. Au lieu de cela\unsure{.|,}
je me suis amusé à repasser en revue un certain nombre de ``couples'' yin-yang.\unsure{\inserted{:}}
Commençant par ceux qui me passaient par la tête, un peu au-bonheur-la-chance,
je me suis ensuite piqué au jeu, et ai terminé par une sorte de ``recensement'' de
tous ceux sur lesquels j'arrivais à mettre la main. J'avais commencé parce que
je m'étais dit que pas mal de ce que j'avais écrit dernièrement risquait fort de
passer entièrement ``par dessus la tête'' d'un lecteur qui ne serait déjà tant \unsure{soi|soit}
peu familier du double aspect yin-yang des choses. Il ne serait peut-être pas
inutile de prendre la peine de donner au moins quelques exemples frappants de
tels couples, en plus de ceux qui s'étaient introduits par la bande ces derniers
% position: left foot mark
\margin{\unsure{-}}
% position: foot pencil
\margin{160}
